%!TEX program = xelatex
\documentclass[12pt, a4paper]{ctexart}

% ===== 页面设置 =====
\usepackage{geometry}
\geometry{left=3cm, right=3cm, top=2.5cm, bottom=2.5cm}

% ===== 行距 =====
\linespread{1.4}

\begin{document}

\begin{center}
  {\zihao{2}\heiti 毛泽东长征诗词三首}
\end{center}

\vspace{1cm}

% ========== 第一首：七律·长征 ==========
\begin{center}
  {\zihao{3}\heiti 七律·长征}\\[0.2cm]
  {\zihao{4}\kaishu 毛泽东\quad 1935年10月}
\end{center}

\vspace{0.3cm}

\begin{center}
  {\zihao{4}\kaishu
  \begin{tabular}{c}
    红军不怕远征难，万水千山只等闲。 \\
    五岭逶迤腾细浪，乌蒙磅礴走泥丸。 \\
    金沙水拍云崖暖，大渡桥横铁索寒。 \\
    更喜岷山千里雪，三军过后尽开颜。
  \end{tabular}}
\end{center}

\vspace{1.2cm}

% ========== 第二首：清平乐·六盘山 ==========
\begin{center}
  {\zihao{3}\heiti 清平乐·六盘山}\\[0.2cm]
  {\zihao{4}\kaishu 毛泽东\quad 1935年10月}
\end{center}

\vspace{0.3cm}

\begin{center}
  {\zihao{4}\kaishu
  \begin{tabular}{c}
    天高云淡，望断南飞雁。 \\
    不到长城非好汉，屈指行程二万。 \\[0.4cm]
    六盘山上高峰，红旗漫卷西风。 \\
    今日长缨在手，何时缚住苍龙？
  \end{tabular}}
\end{center}

\vspace{1.2cm}

% ========== 第三首：忆秦娥·娄山关 ==========
\begin{center}
  {\zihao{3}\heiti 忆秦娥·娄山关}\\[0.2cm]
  {\zihao{4}\kaishu 毛泽东\quad 1935年2月}
\end{center}

\vspace{0.3cm}

\begin{center}
  {\zihao{4}\kaishu
  \begin{tabular}{c}
    西风烈，长空雁叫霜晨月。 \\
    霜晨月，马蹄声碎，喇叭声咽。 \\[0.4cm]
    雄关漫道真如铁，而今迈步从头越。 \\
    从头越，苍山如海，残阳如血。
  \end{tabular}}
\end{center}

\end{document}